\documentclass[10.5pt]{article}
\usepackage[letterpaper, margin=0.6in, top=0.5in, bottom=0.5in]{geometry}
\usepackage{amsmath,amssymb,amsthm}
\usepackage{fancyhdr}
\usepackage{xcolor}
\usepackage{tikz}
\usepackage{enumitem}
\usepackage{tcolorbox}
\usepackage{hyperref}

% Define colors matching instructor board annotations
\definecolor{chalkblue}{RGB}{0, 105, 185}
\definecolor{chalkred}{RGB}{205, 35, 35}
\definecolor{chalkpink}{RGB}{230, 90, 120}
\definecolor{chalkcyan}{RGB}{15, 160, 190}
\definecolor{chalkgreen}{RGB}{20, 130, 50}
\definecolor{boxbg}{RGB}{252, 252, 254}
\definecolor{boxframe}{RGB}{150, 160, 180}

\pagestyle{fancy}
\fancyhf{}
\renewcommand{\headrulewidth}{0pt}
\renewcommand{\footrulewidth}{0pt}

\setlength{\parindent}{0pt}
\setlength{\parskip}{2.5pt}

\begin{document}

% ==============================================================================
% PAGE 1
% ==============================================================================
\noindent \textbf{\large Math 105L} \hfill \textbf{\large Lecture 8-3}\\[1pt]
{\textbf{\Large The Chain Rule}}\\[-5pt]
\rule{\textwidth}{1.2pt}

\vspace{0.03in}

\textbf{\large Past Differentiation\dots}

The differentiation rules we've learned so far:
\begin{itemize}[leftmargin=1.2em, topsep=1pt, itemsep=2pt]
  \item \textbf{Power Rule:} $\frac{d}{dx}[x^n] = n x^{n-1}$
  \item \textbf{Sum Rule:} $\frac{d}{dx}[f(x) + g(x)] = f'(x) + g'(x)$
  \item \textbf{Difference Rule:} $\frac{d}{dx}[f(x) - g(x)] = f'(x) - g'(x)$
  \item \textbf{Product Rule:} $\frac{d}{dx}[f(x)g(x)] = f'(x)g(x) + f(x)g'(x)$
  \item \textbf{Quotient Rule:} $\frac{d}{dx}\left[\frac{f(x)}{g(x)}\right] = \frac{f'(x)g(x) - f(x)g'(x)}{(g(x))^2}$
\end{itemize}

\vspace{0.03in}

\textbf{\large Preliminaries: What We Know So Far}

\textbf{Question: With the tools we have, how do we find $\frac{d}{dx} (x^3 + 1)^3$?}
\begin{itemize}[leftmargin=1.2em, label=$\bullet$, topsep=1pt, itemsep=2pt]
  \item \textbf{Expand algebraically first:}
  \[
    (x^3 + 1)^3 = (x^3 + 1)^2(x^3 + 1) = (x^6 + 2x^3 + 1)(x^3 + 1) = x^9 + 3x^6 + 3x^3 + 1
  \]
  \item \textbf{Differentiate term-by-term via Power Rule:}
  \[
    \frac{d}{dx}(x^3 + 1)^3 = \mathbf{9x^8 + 18x^5 + 9x^2}
  \]
\end{itemize}

\textbf{Question: Does the same idea practically let us find $\frac{d}{dx}(x^3 + 1)^{100}$?}
\begin{itemize}[leftmargin=1.2em, label=$\bullet$, topsep=0pt, itemsep=1pt]
  \item \textbf{Board Answer:} \textit{Not really! Much too long!} We need a rule specifically for function composition.
\end{itemize}

\vspace{0.03in}

\textbf{\large Practical Motivation via Units of Change}

Let $L(E)$ be the life expectancy (years) of a person with $E$ years of education. Let $E(i)$ be the average number of years of education obtained by people with an average annual income of $i$ dollars.

\begin{itemize}[leftmargin=1.2em, topsep=1pt, itemsep=2pt]
  \item \textbf{Describe $L(E(i))$ in practical terms:} Life expectancy (in years) of a person whose parents have an annual income of $\$i$.
  \item \textbf{Units of Derivatives:}
  \[
    \frac{dL}{dE} = \frac{\text{years of life expectancy}}{\text{year of education}}, \qquad \frac{dE}{di} = \frac{\text{years of education}}{\$ \text{ of annual income}}
  \]
  \item \textbf{Chain of Rates:} Multiplying the rates cancels the intermediate unit:
  \[
    \frac{dL}{di} = \frac{dL}{dE} \cdot \frac{dE}{di} = \frac{\text{years of life expectancy}}{\$ \text{ of annual income}}
  \]
\end{itemize}

\vspace{0.03in}

\begin{tcolorbox}[colback=chalkblue!6, colframe=chalkblue, boxrule=1pt, arc=3pt, left=8pt, right=8pt, top=4pt, bottom=4pt]
\textbf{\large The Chain Rule}
\vspace{1pt}

If $f(x)$ and $g(x)$ are differentiable functions, then the composite function is differentiable and:
\[
  \mathbf{\frac{d}{dx}\big[f(g(x))\big] = f'\big(g(x)\big) \cdot g'(x)}
\]
\textbf{Leibniz Notation:} If we set $u = g(x)$ and $y = f(u)$, then $y = f(g(x))$ and:
\[
  \mathbf{\frac{dy}{dx} = \frac{dy}{du} \cdot \frac{du}{dx}}
\]
\end{tcolorbox}

\vspace{0.02in}

\textbf{Example: Differentiating $(x^3 + 1)^3$ via the Chain Rule}
\begin{itemize}[leftmargin=1.2em, label=$\bullet$, topsep=1pt, itemsep=2pt]
  \item Outer function: $f(u) = u^3 \implies f'(u) = 3u^2$
  \item Inner function: $g(x) = x^3 + 1 \implies g'(x) = 3x^2$
  \item Composition: $f'(g(x)) \cdot g'(x) = 3(x^3 + 1)^2 \cdot (3x^2) = 9x^2(x^6 + 2x^3 + 1) = \mathbf{9x^8 + 18x^5 + 9x^2}$ \quad \textit{(Matches!)}
\end{itemize}

\newpage

% ==============================================================================
% PAGE 2
% ==============================================================================
\noindent \textbf{\large Math 105L} \hfill \textbf{\large Lecture 8-3 (Page 2)}\\[1pt]
{\textbf{\Large Proving the Quotient Rule \& Multi-Layer Chain Rules}}\\[-5pt]
\rule{\textwidth}{1.2pt}

\vspace{0.03in}

\textbf{\large Proving the Quotient Rule Using the Chain Rule}

The quotient rule tells us how to differentiate $\frac{f(x)}{g(x)}$. Rather than using difference quotients from scratch, we can prove it naturally by combining the \textbf{Product Rule} and the \textbf{Chain Rule}.

\begin{itemize}[leftmargin=1.2em, topsep=1pt, itemsep=2pt]
  \item \textbf{Rewrite as a Product of Powers:}
  \[
    \frac{f(x)}{g(x)} = f(x) \cdot [g(x)]^{-1}
  \]
  \item \textbf{Step 1: Differentiate the Reciprocal $[g(x)]^{-1}$ via the Chain Rule:}
  \[
    \text{outer: } u^{-1} \implies \frac{d}{du}[u^{-1}] = -u^{-2} = -\frac{1}{u^2}, \qquad \text{inner: } u = g(x) \implies u' = g'(x)
  \]
  \[
    \frac{d}{dx}\big([g(x)]^{-1}\big) = -[g(x)]^{-2} \cdot g'(x) = \mathbf{-\frac{g'(x)}{(g(x))^2}}
  \]
  \item \textbf{Step 2: Differentiate the Product $f(x) \cdot [g(x)]^{-1}$:}
  \begin{align*}
    \frac{d}{dx}\left[f(x)[g(x)]^{-1}\right] &= f'(x)[g(x)]^{-1} + f(x) \frac{d}{dx}\big([g(x)]^{-1}\big) \\
    &= \frac{f'(x)}{g(x)} + f(x)\left(-\frac{g'(x)}{(g(x))^2}\right) \\
    &= \frac{f'(x)g(x)}{(g(x))^2} - \frac{f(x)g'(x)}{(g(x))^2} \\
    &= \mathbf{\frac{f'(x)g(x) - f(x)g'(x)}{(g(x))^2}} \qquad \checkmark
  \end{align*}
\end{itemize}

\vspace{0.04in}

\textbf{\large The Triple Chain Rule (and Higher Compositions)}

What if a function consists of three or more nested layers, such as $f(g(h(x)))$?

\begin{tcolorbox}[colback=boxbg, colframe=boxframe, boxrule=0.8pt, arc=3pt, left=8pt, right=8pt, top=4pt, bottom=4pt]
\textbf{The Triple Chain Rule:}
\[
  \mathbf{\frac{d}{dx}\big[f(g(h(x)))\big] = f'\big(g(h(x))\big) \cdot g'\big(h(x)\big) \cdot h'(x)}
\]
\textbf{General $n$-Layer Composition:} Differentiate from the outermost layer to the innermost layer, evaluating each derivative at the remaining nested inner expression, and multiplying by the derivative of each successive layer.
\end{tcolorbox}

\vspace{0.03in}

\textbf{Example Worked on Blackboard: Differentiating $\mathbf{\big((x^3 + 1)^7 - 3\big)}$}

Identify the structural decomposition across the three layers:
\begin{center}
\renewcommand{\arraystretch}{1.3}
\begin{tabular}{|l|l|l|l|}
\hline
\textbf{Layer} & \textbf{Component Function} & \textbf{Standard Derivative} & \textbf{Evaluated at Inner Argument} \\ \hline
\textbf{Innermost} & $h(x) = x^3 + 1$ & $h'(x) = 3x^2$ & $h'(x) = 3x^2$ \\ \hline
\textbf{Middle} & $g(u) = u^7$ & $g'(u) = 7u^6$ & $g'(h(x)) = 7(x^3 + 1)^6$ \\ \hline
\textbf{Outermost} & $f(v) = v - 3$ & $f'(v) = 1$ & $f'(g(h(x))) = 1$ \\ \hline
\end{tabular}
\end{center}

\vspace{0.02in}

Applying the Triple Chain Rule:
\begin{align*}
  \frac{d}{dx}\Big[\big((x^3 + 1)^7 - 3\big)\Big] &= f'\big(g(h(x))\big) \cdot g'\big(h(x)\big) \cdot h'(x) \\
  &= 1 \cdot \left[7(x^3 + 1)^6\right] \cdot \left(3x^2\right) \\
  &= \mathbf{21x^2 (x^3 + 1)^6}
\end{align*}

\end{document}

